\begin{textsample}{NEG Dim 1 – 2009\_c4 – Score -9.00 – 2009\_c4/enade\_2009\_com\_soc\_edit\_bd87747c88.txt}  \label{ex:f1_neg_023}
Capítulo 4 Percepção sobre a Prova As percepções dos estudantes sobre a prova aplicada no ENADE/2009 na área de COMUNICAÇÃO SOCIAL – EDITORAÇÃO foram mensuradas por meio de nove itens que avaliaram desde o grau de dificuldade da prova até o tempo gasto para concluí -la.
As questões foram analisadas separando -se concluintes e ingressantes, e as percepções sobre a prova foram relacionadas com o desempenho dos estudantes e com a Região Geográfica de \textbf{localização} do curso6.
O desempenho dos estudantes dentro de \textbf{cada} área/Subárea foi classificado em quartos ( 1/4 ) para os concluintes e para os ingressantes em separado ( ver Figura 1 para o \textbf{esquema} de representação gráfica destas grandezas ).
As notas finais foram colocadas em \textbf{ordem} ascendente e no primeiro quarto foram colocadas as 25 \% notas mais baixas, limitadas superiormente pelo percentil 25 que vem a ser o primeiro quartil.
No segundo quarto, as 25 \% seguintes, limitadas superiormente pelo percentil 50 que vem a ser a mediana.
No terceiro quarto, as 25 \% notas após as encontradas no segundo quarto, ou seja, as primeiras 25 \% acima da mediana, limitadas superiormente pelo percentil 75 que vem a ser o terceiro quartil.
No último quarto foram alocadas as 25 \% maiores notas.
Potencialmente, poderia existir uma relação entre a nota do estudante e a sua percepção sobre a prova.
Calculando as médias das respostas segundo os quartos de notas, seria possível reconhecer a existência ( ou não ) de algum gradiente nas respostas, por exemplo, uma melhor nota correspondendo a uma percepção menos negativa.
Este gradiente seria traduzido, por exemplo no gráfico 4.9a, pela diminuição das áreas amarela, laranja e vermelha ( “ Apenas cerca de metade ”, “ Poucos ” e “ Não, nenhum ” ) com o aumento do desempenho ( indo do primeiro ao último quarto ). 6 Curso, neste relatório, refere -se à unidade de análise para o Conceito ENADE e é caracterizado pela combinação de IES, Área/Subárea ( quando pertinente ) e município de habilitação.
A medida tradicional de “ grau de discriminação ” comparou os resultados dos estudantes do primeiro quarto ( 25 \% piores estudantes ) aos do último quarto ( 25 \% melhores ).
Tal comparação, no gráfico disponibilizado, foi idêntica à comparação da primeira barra ( primeiro quarto ) com a última barra ( último quarto ) de \textbf{cada} grupo ingressante ou concluinte nos gráficos “ segundo quartos de nota final ” ( gráficos com números ímpares neste capítulo ).
Por outro lado, a disponibilização das informações segundo os quartos de nota permitiu verificar o comportamento de dois outros grupos intermediários que na análise somente dos extremos tem que ser suposto monotônico.
A seguir, serão apresentados os principais resultados relativos aos nove itens avaliados. 4.1 Grau de dificuldade da prova 4.1.1 Formação Geral Os Gráficos 4.1a, 4.1b, 4.1c apresentam a distribuição das respostas à questão “ Qual o grau de dificuldade desta prova na parte de Formação Geral? ” ; segundo quartos de nota final e grupo de estudantes concluintes e ingressantes que participaram do ENADE/2009 na área de COMUNICAÇÃO SOCIAL – EDITORAÇÃO.
No total de concluintes, 29,5 \% declararam que a prova na parte da Formação Geral foi Fácil/Muito fácil e um pouco abaixo os ingressantes, 23,4 \%.
A proporção que declarou a parte Difícil/Muito difícil foi menor entre os concluintes ( 9,9 \% ) do que entre os ingressantes ( 13,2 \% ).
Não se observou uma variação monotônica significativa desta proporção com o aumento no desempenho, nem entre os concluintes, nem entre os ingressantes.
Entre os concluintes, 60,6 \% classificaram esta parte com um nível Médio de dificuldade.
A proporção entre os ingressantes foi ligeiramente maior, 63,4 \%.
Os Gráficos 4.2a e 4.2b apresentam a distribuição das respostas à questão “ Qual o grau de dificuldade desta prova na parte de Formação Geral? ”, segundo a Região onde o curso \textbf{está} situado para os estudantes concluintes e ingressantes que participaram do ENADE/2009 na área de COMUNICAÇÃO SOCIAL – EDITORAÇÃO.
Cumpre \textbf{lembrar} que a Região Norte não apresentou nenhum estudante concluinte e a Região Centro-Oeste não apresentou nenhum estudante, quer seja concluinte ou ingressante.
Na Região Sul, nenhum estudante concluinte declarou a prova como Muito difícil.
Entre os concluintes, a maior proporção de estudantes que considerou a parte de Formação Geral da prova como de dificuldade média entre todas as Regiões brasileiras, foram os da Região Sul com 80,0 \%, entre os ingressantes, foram também os estudantes da Região Sul com 72,4 \%.
A Região com as menores proporções desta resposta foram respectivamente para os concluintes e ingressantes, a Nordeste ( 42,9 \% ) e a Norte ( 43,3 \% ). 4.1.2 Componente Específico Os Gráficos 4.3a, 4.3b e 4.3c apresentam a distribuição das respostas à questão “ Qual o grau de dificuldade desta prova na parte do Componente Específico? ”, segundo quartos de nota final e grupo de estudantes concluintes e ingressantes que participaram do ENADE/2009 na área de COMUNICAÇÃO SOCIAL – EDITORAÇÃO.
Em geral, para as desagregações consideradas, os estudantes classificaram o Componente Específico como mais difícil que a Formação Geral.
No grupo de concluintes, 10,7 \% declarou que a prova na parte do Componente Específico foi Difícil/Muito difícil, sendo um valor muito maior para os ingressantes, 19,6 \%.
A proporção que declarou a parte Fácil/Muito fácil foi maior entre os concluintes ( 25,7 \% ) do que entre os ingressantes ( 16,5 \% ).
Entre os concluintes, 63,6 \% classificaram esta parte com um nível Médio de dificuldade.
A proporção entre os ingressantes foi ligeiramente maior, 64,0 \%.
Os Gráficos 4.4a e 4.4b apresentam a distribuição das respostas à questão “ Qual o grau de dificuldade desta prova na parte do Componente Específico? ”, segundo a Região onde o curso \textbf{está} situado para os estudantes concluintes e ingressantes que participaram do ENADE/2009 na área de COMUNICAÇÃO SOCIAL – EDITORAÇÃO.
Cumpre \textbf{lembrar} aqui também que a Região Norte não apresentou nenhum estudante concluinte e a Centro-Oeste nem concluinte e nem ingressante.
Entre os concluintes, a maior proporção de estudantes que considerou a parte do Componente Específico da prova como de dificuldade média entre todas as Regiões brasileiras, foi na Região Sudeste com 68,9 \%, e entre os ingressantes foi também na Região Sudeste com 68,5 \%. 4.2 Extensão da prova em relação ao tempo total Os Gráficos 4.5a, 4.5b e 4.5c apresentam a distribuição das respostas à questão sobre a extensão da prova em relação ao tempo total segundo quartos de nota final e grupo de estudantes concluintes e ingressantes que participaram do ENADE/2009 na área de COMUNICAÇÃO SOCIAL – EDITORAÇÃO.
Tanto para concluintes quanto para ingressantes, a maioria considerou a extensão da prova como Adequada em relação ao tempo destinado à sua resolução, respectivamente 54,0 \% e 61,6 \%.
A proporção dos que consideraram a extensão da prova Muito curta foi maior entre os concluintes ( 2,2 \% ) do que entre os ingressantes ( 0,9 \% ).
A proporção dos que consideraram a extensão da prova Curta/Muito curta vis-à-vis ao tempo total foi maior entre os concluintes ( 7,2 \% ) do que entre os ingressantes ( 6,9 \% ).
Entre os concluintes, nenhum estudante com desempenho no 2º grupo declarou que a extensão da prova teria sido Muito curta com relação ao tempo total.
Entre os ingressantes, esta declaração não aconteceu no 2º e no 3º quarto.
A maior proporção de concluintes declarando que a prova foi Curta/Muito curta foi entre os estudantes de pior desempenho, 8,9 \%.
A proporção dos que consideraram a extensão da prova Muito longa foi próxima entre os concluintes ( 9,4 \% ) e ingressantes ( 12,1 \% ).
A proporção dos que consideraram a extensão da prova Longa foi maior entre os concluintes ( 29,5 \% ) do que entre os ingressantes ( 19,4 \% ).
Os Gráficos 4.6a e 4.6b apresentam a distribuição das respostas à questão sobre a extensão da prova em relação ao tempo total, segundo a Região onde o curso \textbf{está} situado para os estudantes concluintes e ingressantes que participaram do ENADE/2009 na área de COMUNICAÇÃO SOCIAL – EDITORAÇÃO.
Em todas as Regiões brasileiras, a maioria dos estudantes, tanto concluintes, como ingressantes, percebeu a prova como adequada.
Não houve estudantes da Região Centro-Oeste.
Os concluintes das Regiões Sudeste ( 57,1 \% ) e Sul ( 53,3 \% ) se situaram acima da média ( 54,0 \% ) para esta resposta.
Os ingressantes das Regiões Norte ( 66,7 \% ) e Nordeste ( 64,1 \% ) se situaram acima da média ( 61,6 \% ) para esta resposta. 4.3 Compreensão dos Enunciados das Questões 4.3.1 Formação Geral Os Gráficos 4.7a, 4.7b e 4.7c apresentam a distribuição das respostas à questão “ Os enunciados das questões da prova na parte de Formação Geral \textbf{estavam} claros e objetivos? ”, segundo quartos de nota final e grupo de estudantes concluintes e ingressantes que participaram do ENADE/2009 na área de COMUNICAÇÃO SOCIAL – EDITORAÇÃO.
Com relação aos enunciados das questões da parte de Formação Geral, as opiniões foram positivas : a grande maioria dos concluintes ( 73,5 \% ) e ingressantes ( 83,3 \% ) consideram claros e objetivos os enunciados de todas ou da maioria das questões.
Entre os concluintes, esta proporção cresce com o desempenho : 58,6 \% para o grupo com o pior desempenho e 78,9 \% para o de melhor.
A proporção dos que consideram a combinação de “ Poucos ” e “ Não, nenhum ” foi maior entre os concluintes ( 11,4 \% ) do que entre os ingressantes ( 6,4 \% ).
Os Gráficos 4.8a e 4.8b apresentam a distribuição das respostas à questão “ Os enunciados das questões da prova na parte de Formação Geral \textbf{estavam} claros e objetivos? ”, segundo a Região onde o curso \textbf{está} situado para os estudantes concluintes e ingressantes que participaram do ENADE/2009 na área de COMUNICAÇÃO SOCIAL – EDITORAÇÃO.
Em termos regionais, as proporções de ingressantes e de concluintes que consideraram claros e objetivos os enunciados de todas ou da maioria das questões foram respectivamente, 83,3 \% e 73,6 \%.
Entre os concluintes, os valores variaram de 66,7 \% no Sul a 75,0 \% no Nordeste.
Já entre os ingressantes estas proporções variaram de 75,0 \% no Sul a 86,7 \% no Sudeste.
Com relação às respostas “ poucos ” e “ não, nenhum ”, as maiores proporções para concluintes ocorreram na Região Sul ( 13,3 \% ) e para ingressantes na Região Nordeste ( 7,5 \% ). 4.3.2 Componente Específico Os Gráficos 4.9a, 4.9b e 4.9c apresentam a distribuição das respostas à questão “ Os enunciados das questões da prova na parte de formação específica \textbf{estavam} claros e objetivos? ”, segundo quartos de nota final e grupo de estudantes concluintes e ingressantes que participaram do ENADE/2009 na área de COMUNICAÇÃO SOCIAL – EDITORAÇÃO.
Com relação aos enunciados das questões da parte do Componente Específico, as opiniões foram também positivas : a grande maioria dos concluintes ( 75,3 \% ) e ingressantes ( 80,9 \% ) consideraram claros e objetivos os enunciados de todas ou da maioria das questões.
Já a combinação de todas ou da maioria das questões foi crescente entre os concluintes com o desempenho, indo de 63,4 \% para os com pior desempenho a 86,5 \% para o último quarto.
Entre os ingressantes, esta proporção cresce com o desempenho indo de 75,0 \% para os com pior desempenho a 88,8 \% para os com melhor desempenho.
O número de estudantes que declararam que “ não, nenhum ” dos enunciados das questões da parte do Componente Específico era claro e objetivo, foi maior do que na parte de Formação Geral ( 2,1 \% entre os concluintes e ingressantes ) : 3,6 \% entre os concluintes e 2,7 \% entre os ingressantes.
Os que responderam “ poucos ”, são 6,5 \% entre os concluintes e 4,4 \% entre os ingressantes.
A proporção de estudantes que respondeu que cerca de metade dos enunciados da parte do Componente Específico era claro e objetivo decresceu com o desempenho para concluintes.
Os Gráficos 4.10a e 4.10b apresentam a distribuição das respostas à questão “ Os enunciados das questões da prova na parte de formação específica \textbf{estavam} claros e objetivos? ”, segundo a Região onde o curso \textbf{está} situado para os estudantes concluintes e ingressantes que participaram do ENADE/2009 na área de COMUNICAÇÃO SOCIAL – EDITORAÇÃO.
Em termos regionais, as proporções de ingressantes e de concluintes que consideraram claros e objetivos os enunciados de todas ou da maioria das questões variaram um pouco menos entre os ingressantes ( no máximo 16 pontos percentuais ) do que entre os concluintes ( no máximo 17 pontos percentuais ). 4.4 Suficiência das informações/instruções fornecidas Os Gráficos 4.11a, 4.11b e 4.11c apresentam a distribuição das respostas à questão “ As informações/instruções fornecidas para a resolução das questões foram suficientes para resolvê -las? ”, segundo quartos de nota final e grupo de estudantes concluintes e ingressantes que participaram do ENADE/2009 na área de COMUNICAÇÃO SOCIAL – EDITORAÇÃO.
Quando \textbf{perguntados} sobre as informações/instruções fornecidas para a resolução das questões, 86,2 \% dos ingressantes e 86,6 \% dos concluintes de todo o Brasil responderam que estas eram ou “ excessivas ” ou “ suficientes em todas ” ou “ na maioria das questões ”, demonstrando uma avaliação positiva da prova.
A resposta “ não, em nenhuma delas ” apareceu somente no 1º quarto para concluintes, e apareceu em todos os quartos para ingressantes independente do desempenho.
Os Gráficos 4.12a e 4.12b apresentam a distribuição das respostas à questão “ As informações/instruções fornecidas para a resolução das questões foram suficientes para resolvê -las? ”, segundo a Região onde o curso \textbf{está} situado para os estudantes concluintes e ingressantes que participaram do ENADE/2009 na área de COMUNICAÇÃO SOCIAL – EDITORAÇÃO.
Em termos regionais, observa -se que a proporção de estudantes concluintes que perceberam como “ excessivas ” ou “ suficientes ” as informações/instruções fornecidas para a resolução de todas ou da “ maioria ” das questões superaram a média nacional ( 86,6 \% ) na Região Sudeste ( 88,7 \% ).
Já entre os ingressantes, a média nacional foi 86,3 \% e ficaram acima deste valor as Regiões Sul ( 93,1 \% ) e Sudeste ( 88,7 \% ).
Cabe salientar que a Região Centro-Oeste não tem estudantes e que a Região Norte não tem estudantes concluintes. 4.5 Dificuldade encontrada ao responder à prova Os Gráficos 4.13a, 4.13b e 4.13c apresentam a distribuição das respostas à questão “ Você se deparou com alguma dificuldade ao responder à prova?
Qual? ”, segundo quartos de nota final e grupo de estudantes concluintes e ingressantes que participaram do ENADE/2009 na área de COMUNICAÇÃO SOCIAL – EDITORAÇÃO. \textbf{Perguntados} sobre as dificuldades com as quais se depararam ao responder à prova, 31,5 \% dos ingressantes apontaram como principal dificuldade o desconhecimento do conteúdo, já os concluintes apontaram como principal dificuldade a forma diferente de abordagem, 32,4 \%.
Na verdade, para a quase totalidade das áreas, ingressantes apresentaram uma proporção de estudantes bem maior do que entre os concluintes declarando como principal dificuldade o desconhecimento da matéria.
O desconhecimento do conteúdo foi apontado somente por 13,7 \% dos concluintes.
A forma diferente de abordagem do conteúdo foi apontada por 26,7 \% dos ingressantes, enquanto que a falta de motivação foi a principal dificuldade para 19,4 \% dos ingressantes.
Esta dificuldade apresentou valores no entorno de 27 \% para os três últimos grupos de desempenho de concluintes.
Entre os ingressantes, esses valores são mais homogêneos em \textbf{torno} de 21 \%.
Em termos nacionais, a proporção de ingressantes e concluintes que declararam que não tiveram qualquer dificuldade para responder à prova foi, respectivamente, 18,1 \% e 14,4 \%.
Os Gráficos 4.14a e 4.14b apresentam a distribuição das respostas à questão “ Você se deparou com alguma dificuldade ao responder à prova? ”, segundo a Região onde o curso \textbf{está} situado para os estudantes concluintes e ingressantes que participaram do ENADE/2009 na área de COMUNICAÇÃO SOCIAL – EDITORAÇÃO.
Em relação às Regiões brasileiras, o percentual de ingressantes que apontaram o desconhecimento do conteúdo como principal dificuldade ao responder à prova variou de 20,7 \% na Região Sul a 35,1 \% na Sudeste.
Entre os concluintes, tal percentual variou de 7,7 \% na Região Sudeste a 42,9 \% na Nordeste.
A forma diferente de abordagem de conteúdo foi citada como maior dificuldade por 23,8 \% dos concluintes na Região Nordeste chegando a 57,1 \% na Região Sul.
Já entre os ingressantes os valores foram de 15,4 \% na Região Nordeste a 50,0 \% na Região Norte.
O percentual de concluintes que citaram a falta de motivação como principal dificuldade variou de 21,4 \% na Região Sul a 31,7 \% na Região Sudeste.
Entre os ingressantes, as estatísticas correspondentes foram de 3,3 \% na Região Norte a 35,9 \%, na Região Nordeste.
Os que declararam que não tiveram qualquer tipo de dificuldade variaram, entre os concluintes, de 4,8 \% na Região Nordeste a 17,3 \% na Região Sudeste.
Entre os ingressantes, este percentual variou de 13,3 \% na Região Norte a 20,7 \% na Região Sul. 4.6 Percepção sobre os conteúdos das questões objetivas da prova Os Gráficos 4.15a, 4.15b e 4.15c apresentam a distribuição das respostas à questão “ Considerando apenas as questões objetivas da prova, você percebeu que :... ”, segundo quartos de nota final e grupo de estudantes concluintes e ingressantes que participaram do ENADE/2009 na área de COMUNICAÇÃO SOCIAL – EDITORAÇÃO.
As respostas possíveis para esta questão \textbf{estavam} relacionadas à \textbf{cobertura} do conteúdo das questões objetivas durante o curso.
A proporção de ingressantes ( 40,4 \% ) que declarou que não estudou ainda a maioria desses conteúdos foi bem maior do que entre os concluintes ( 9,2 \% ).
Esta proporção foi decrescente para os ingressantes como função do desempenho.
Por outro lado, a proporção de estudantes que declarou que estudou e aprendeu todos ou a maioria desses conteúdos, entre os ingressantes foi 46,9 \% e entre os concluintes, 77,5 \%.
Esta proporção foi crescente como função do desempenho, tanto para os concluintes como para os ingressantes.
A proporção de estudantes que declarou ter estudado a maioria desses conteúdos, mas não os ter aprendido é maior entre os concluintes ( 5,6 \% ) do que entre os ingressantes ( 4,3 \% ).
Os Gráficos 4.16a e 4.16b apresentam a distribuição das respostas à questão “ Considerando apenas as questões objetivas da prova, você percebeu que :... ”, para os estudantes concluintes e ingressantes que participaram do ENADE/2009 na área de COMUNICAÇÃO SOCIAL – EDITORAÇÃO, segundo região onde o curso \textbf{está} situado.
Em relação às Regiões brasileiras, concluintes que afirmaram “ não ter estudado ainda a maioria dos conteúdos ” só ocorreram nas Regiões Nordeste ( 23,8 \% ), Sudeste ( 6,6 \% ) e Sul ( 6,7 \% ).
Entre os ingressantes, os percentuais variaram de 33,1 \% na Região Sudeste a 51,7 \%, na Norte.
Em todas as Regiões, a maioria dos concluintes afirmou"ter estudado e aprendido todos ou muitos dos conteúdos ”, com percentuais que variaram de 61,9 \% na Região Nordeste a 86,6 \%, na Sul.
Para os ingressantes, tal percentual variou de 34,5 \%, na Região Norte, a 53,4 \%, na Sudeste. 4.7 Tempo gasto para concluir a prova Os Gráficos 4.17a, 4.17b e 4.17c apresentam a distribuição das respostas à questão “ Qual foi o tempo gasto para concluir a prova? ”, segundo quartos de nota final e grupo de estudantes concluintes e ingressantes que participaram do ENADE/2009 na área de COMUNICAÇÃO SOCIAL – EDITORAÇÃO.
Quando \textbf{perguntados} quanto ao tempo de conclusão da prova, a maioria dos concluintes respondeu entre uma e duas horas ( 35,5 \% ), já os ingressantes responderam entre duas e três horas ( 41,5 \% ).
Para a segunda maior incidência de respostas ocorre uma inversão de posição, pois a opção entre uma e duas horas é o segundo maior percentual para os ingressantes ( 39,7 \% ) e a opção em segundo lugar para os concluintes ( 34,0 \% ) é entre duas e três horas.
A proporção daqueles que declararam ter usado as quatro horas, mas não terem conseguido terminar foi bem baixa : 1,7 \% entre os ingressantes e 1,4 \% entre os concluintes.
Entre os concluintes, esta resposta ficou concentrada exclusivamente entre os estudantes do primeiro quarto, aqueles com pior desempenho.
Os Gráficos 4.18a e 4.18b apresentam a distribuição das respostas à questão “ Qual foi o tempo gasto por você para concluir a prova? ”, segundo a Região onde o curso \textbf{está} situado para os estudantes concluintes e ingressantes que participaram do ENADE/2009 na área de COMUNICAÇÃO SOCIAL – EDITORAÇÃO.
A proporção de concluintes que gastaram entre uma e duas horas foi maior na Região Nordeste ( 42,9 \% ) e menor na Região Sul ( 20,0 \% ).
Para ingressantes, o maior percentual para este tempo ocorreu na Região Nordeste ( 51,3 \% ) e a menor na Região Norte ( 23,3 \% ).

% matched lemmas: cada, cobertura, esquema, estar, lembrar, localização, ordem, perguntar, torno
\end{textsample}
